\chapter{Learning Theory — Halving}
%%% generated by the blueprint pipeline from TCSlib.LearningTheory.Halving
\section{Overview}

This module formalises the Halving Algorithm for binary classification in the
realizable online-learning setting.  Given a finite version space $V$ of hypotheses
and a sequence of labelled examples generated by an unknown target hypothesis, the
algorithm predicts by majority vote and eliminates every consistent-minority on each
mistake.  The central result is that the number of mistakes is at most
$\lfloor\log_2 |V|\rfloor$.

%%%%%%%%%%%%%%%%%%%%%%%%%%%%%%%%%%%%%%%%%%%%%%%%%%%%%%%%%%%%%%%%%%%%%%%%%%%%%%%
\section{Declarations}

\begin{definition}[Vote set for a label]
\lean{Halving.voteFor}
\label{Halving.voteFor}
\leanok
Given an evaluation map $\mathtt{eval} : \mathit{Hyp} \to X \to \mathsf{Bool}$, a
version space $V \subseteq \mathit{Hyp}$, an input $x \in X$, and a label
$y \in \mathsf{Bool}$, $\mathtt{voteFor}(\mathtt{eval}, V, x, y)$ is the subset of $V$
consisting of all hypotheses that predict label $y$ on input $x$.
\end{definition}

\begin{definition}[Halving prediction]
\lean{Halving.predict}
\label{Halving.predict}
\leanok
\uses{Halving.voteFor}
The Halving Algorithm predicts on input $x$ by taking a majority vote of the current
version space $V$: it returns $\mathsf{true}$ whenever the number of hypotheses in $V$
that predict $\mathsf{true}$ is at least as large as the number that predict
$\mathsf{false}$ (ties are broken in favour of $\mathsf{true}$).
\end{definition}

\begin{definition}[Version-space update]
\lean{Halving.update}
\label{Halving.update}
\leanok
\uses{Halving.voteFor}
After observing that the correct label for input $x$ is $y$, the version space is
updated to $\mathtt{update}(\mathtt{eval}, V, x, y) = \mathtt{voteFor}(\mathtt{eval},
V, x, y)$, i.e.\ only the hypotheses that correctly predict $y$ on $x$ are retained.
\end{definition}

\begin{lemma}[Vote counts partition the version space]
\lean{Halving.voteFor_false_add_true}
\label{Halving.voteFor_false_add_true}
\leanok
\uses{Halving.voteFor}
\difficulty{3}
Let $\mathrm{eval}$ be a Boolean evaluation map assigning to each hypothesis and each
input a truth value, let $V$ be a finite version space of hypotheses, and let $x$ be an
input. Every hypothesis in $V$ predicts either $\mathsf{false}$ or $\mathsf{true}$ on
$x$, so the number of hypotheses in $V$ predicting $\mathsf{false}$ on $x$ and the
number predicting $\mathsf{true}$ on $x$ sum to the total number of hypotheses in $V$:
\[
  \abs{\{\,h \in V : \mathrm{eval}(h, x) = \mathsf{false}\,\}} +
  \abs{\{\,h \in V : \mathrm{eval}(h, x) = \mathsf{true}\,\}} = \abs{V}.
\]
\end{lemma}

\begin{lemma}[The updated version space is a subset]
\lean{Halving.update_subset}
\label{Halving.update_subset}
\leanok
\uses{Halving.update, Halving.voteFor}
\difficulty{2}
Fix an evaluation map $\mathtt{eval} : \mathit{Hyp} \to X \to \mathsf{Bool}$, a finite
version space $V \subseteq \mathit{Hyp}$, an input $x \in X$, and a label $y \in
\mathsf{Bool}$. Then the updated version space is contained in the original one:
$\mathtt{update}(\mathtt{eval}, V, x, y) \subseteq V$.
\end{lemma}

\begin{lemma}[Target survives its own update]
\lean{Halving.target_mem_update}
\label{Halving.target_mem_update}
\leanok
\uses{Halving.update, Halving.voteFor}
\difficulty{1}
Let $\mathtt{eval} : \mathit{Hyp} \to X \to \mathsf{Bool}$ be an evaluation map, let $V
\subseteq \mathit{Hyp}$ be a version space, fix an input $x \in X$, and let
$\mathtt{target} \in V$ be a hypothesis. Then $\mathtt{target}$ remains in the version
space obtained by updating $V$ on $x$ with the label that $\mathtt{target}$ itself
assigns, that is, $\mathtt{target} \in \mathtt{update}(\mathtt{eval}, V, x,
\mathtt{eval}\,\mathtt{target}\,x)$.
\end{lemma}

\begin{lemma}[Mistake bound for the Halving algorithm]
\lean{Halving.mistake_halves}
\label{Halving.mistake_halves}
\leanok
\uses{Halving.predict, Halving.update, Halving.voteFor, Halving.voteFor_false_add_true}
\difficulty{4}
Let $\mathtt{eval} : \mathit{Hyp} \to X \to \mathsf{Bool}$ be an evaluation map, let $V
\subseteq \mathit{Hyp}$ be a finite version space, let $x \in X$ be an input, and let $y
\in \mathsf{Bool}$ be its true label. If the Halving prediction on $x$ is wrong, that is
$\mathtt{predict}(\mathtt{eval}, V, x) \ne y$, then the updated version space contains
at most half of the hypotheses:
\[
  2 \cdot \abs{\mathtt{update}(\mathtt{eval}, V, x, y)} \le \abs{V}.
\]
\end{lemma}

\begin{lemma}[Logarithm increment under doubling]
\lean{Halving.log_succ_le_log_of_double_le}
\label{Halving.log_succ_le_log_of_double_le}
\leanok
\difficulty{3}
For natural numbers $a$ and $b$ with $b > 0$ and $2b \le a$,
\[
  \lfloor \log_2 b \rfloor + 1 \;\le\; \lfloor \log_2 a \rfloor,
\]
where $\lfloor \log_2 \cdot \rfloor$ denotes the integer (floor) base-$2$ logarithm.
\end{lemma}

\begin{definition}[Mistake count over a sequence]
\lean{Halving.mistakes}
\label{Halving.mistakes}
\leanok
\uses{Halving.predict, Halving.update}
$\mathtt{mistakes}(\mathtt{eval}, \mathtt{target}, V, xs)$ is the total number of
prediction errors made by the Halving Algorithm when it processes the input list $xs$,
using the version space $V$, with the true label at each step supplied by
$\mathtt{target}$.  It is defined recursively: the empty list yields zero mistakes, and
for a list $x :: xs$ the count is $1$ if the algorithm mispredicts on $x$ (else $0$),
plus the mistakes on the remainder using the updated version space.
\end{definition}

\begin{theorem}[Mistake bound for the Halving Algorithm]
\lean{Halving.mistakes_bound}
\label{Halving.mistakes_bound}
\leanok
\uses{Halving.log_succ_le_log_of_double_le, Halving.mistake_halves, Halving.mistakes, Halving.predict, Halving.target_mem_update, Halving.update, Halving.update_subset}
\difficulty{5}
Let $\mathtt{eval} : \mathit{Hyp} \to X \to \mathsf{Bool}$ be an evaluation map, let $V$
be a finite version space of hypotheses, and fix a target hypothesis $\mathtt{target}
\in V$ whose predictions supply the true label at each step. Then for every input
sequence $xs$, the total number of mistakes made by the Halving Algorithm — which
predicts by majority vote over the current version space and, after each input, retains
only those hypotheses that agreed with the observed label — satisfies
\[
\mathtt{mistakes}(\mathtt{eval}, \mathtt{target}, V, xs) \;\le\; \lfloor \log_2 \abs{V}
\rfloor.
\]
\end{theorem}

\begin{theorem}[Halving mistake bound]
\lean{Halving.halving_bound}
\label{Halving.halving_bound}
\leanok
\uses{Halving.mistakes, Halving.mistakes_bound}
\difficulty{1}
Let $\mathtt{eval} : \mathit{Hyp} \to X \to \mathsf{Bool}$ assign a Boolean label to
each input under each hypothesis, let $H$ be a finite class of hypotheses, and fix a
target hypothesis $\mathtt{target} \in H$ that supplies the correct label
$\mathtt{eval}(\mathtt{target}, x)$ for every input $x$. Then, on any finite sequence
$xs$ of inputs, the number of mistakes made by the Halving Algorithm run from the
initial version space $H$ is at most $\lfloor \log_2 \abs{H} \rfloor$:
\[
\mathtt{mistakes}(\mathtt{eval}, \mathtt{target}, H, xs) \;\le\; \lfloor \log_2 \abs{H}
\rfloor.
\]
\end{theorem}
